% Generated draft: requires mathematical and source review.
\documentclass[11pt]{article}
\usepackage[margin=25mm]{geometry}
\usepackage{fontspec}
\setmainfont{Latin Modern Roman}
\usepackage{amsmath,amssymb,mathtools,unicode-math}
\setmathfont{Latin Modern Math}
\usepackage{enumitem,needspace,xurl,hyperref}
\setlength{\parindent}{0pt}
\setlength{\parskip}{5pt}
\setlength{\emergencystretch}{3em}
\newcommand{\N}{\mathbb{N}}
\newcommand{\Z}{\mathbb{Z}}
\newcommand{\Q}{\mathbb{Q}}
\newcommand{\R}{\mathbb{R}}
\newcommand{\C}{\mathbb{C}}
\newcommand{\F}{\mathbb{F}}
\newcommand{\A}{\mathbb{A}}
\newcommand{\PP}{\mathbb{P}}
\newcommand{\E}{\mathbb{E}}
\newcommand{\eps}{\varepsilon}
\newcommand{\dd}{\,\mathrm{d}}
\begin{document}
\section*{UnsolvedMath ID: 700957}
\textbf{Source problem code:} catalogue entry
\subsection{Definitions and mathematical statement}
Let \(A\) be an associative ring with identity whose additive group \((A,+)\) is torsion-free: for every positive integer \(n\) and every \(a\in A\), \(na=0\) implies \(a=0\). Let \(F\) be a free group, and let \(F^A\) denote its tensor product by \(A\), in the group-theoretic sense of the cited work of A. G. Myasnikov and V. N. Remeslennikov. The source calls \(F^A\) the free exponential group over \(A\) and reports a construction using free products with amalgamation. This notation does not denote an ordinary tensor product of modules. The defining axioms and universal property of this construction are not reproduced in the supplied record, so the questions below retain that source-defined meaning. No rank restriction on \(F\) is stated.

The supplied record retains parts (b), (d), (e), and (f) in the main problem and places parts (a) and (c) in its solved archive. The following preserves all six questions without independently verifying their status.

\begin{description}
\item[(b)] Is \(F^A\) a linear group? Under the usual field-based meaning of linearity, the question is whether, for each \(A\) and \(F\) as above, there exist a field \(K\), a positive integer \(m\), and an injective group homomorphism
\[
F^A\longrightarrow \operatorname{GL}_m(K).
\]
The source does not specify the coefficient field, its characteristic, or whether a broader ring-based meaning of linearity is intended.

The source reports an affirmative answer when the additive cyclic subgroup \(\langle 1_A\rangle=\mathbb Z1_A\) is pure in \((A,+)\). Here purity means
\[
(\mathbb Z1_A)\cap nA=n(\mathbb Z1_A)
\qquad\text{for every positive integer }n,
\]
where \(nA=\{na:a\in A\}\).

\item[(d)] Is the universal theory of \(F^A\) decidable? For a specified first-order language \(L\), its universal theory is the set of all sentences
\[
\forall x_1\cdots\forall x_r\,\varphi(x_1,\ldots,x_r)
\]
that hold in \(F^A\), where \(\varphi\) is quantifier-free in \(L\). Decidability asks for an algorithm deciding membership in this set. The supplied source does not specify whether \(L\) is the ordinary group language or a language containing exponentiation operations indexed by \(A\). It also does not specify an effective presentation of \(A\), any rank restriction, or whether a uniform algorithm as \(A\) and \(F\) vary is requested. These choices remain unresolved.

\item[(e), G. Baumslag] Can free \(A\)-groups be characterized by a length function? The source does not specify the codomain or axioms of the length function, the class of groups within which the characterization is sought, or the proposed necessary and sufficient conditions. Consequently, this part remains a request for a characterization with those parameters unresolved; no particular length-function criterion is asserted here.

\item[(f), G. Baumslag] Does a free \(\mathbb Q\)-group admit a free action on some \(\Lambda\)-tree? More explicitly, for a free group \(F\), does there exist an ordered abelian group \(\Lambda\), a \(\Lambda\)-tree \(T\) in the sense of the cited work of R. Alperin and H. Bass, and an isometric action of \(F^{\mathbb Q}\) on \(T\) that is free? The fixed-point condition for freeness is
\[
\forall g\in F^{\mathbb Q}\setminus\{1\}\ \forall p\in T,
\qquad g\cdot p\ne p.
\]
The source refers to Alperin--Bass for the tree terminology and does not separately state its convention concerning inversions; the intended action convention requires review against that reference.

\item[(a), G. Baumslag; archived)] Is \(F^A\) residually torsion-free nilpotent? Under the usual interpretation of the source wording ``residually nilpotent torsion-free'', this asks whether, for every \(g\in F^A\setminus\{1\}\), there are a torsion-free nilpotent group \(N\) and a homomorphism \(q:F^A\to N\) such that \(q(g)\ne1\). A group is torsion-free if no nonidentity element has finite order. A group \(N\) is nilpotent if its lower central series, defined by \(\gamma_1(N)=N\) and \(\gamma_{j+1}(N)=[\gamma_j(N),N]\), reaches the identity subgroup after finitely many steps. The supplied archive reports an affirmative answer attributed to A. Jaikin-Zapirain (2024). That report is supplied source information, not an independently verified theorem here. The precise interpretation of the original wording should be checked against the cited article.

\item[(c), G. Baumslag; archived)] Is the Magnus homomorphism of \(F^{\mathbb Q}\) into the group of power series over \(\mathbb Q\) faithful? Writing \(\mu\) for the homomorphism intended in the source, the question is
\[
\ker\mu=\{1\}.
\]
The supplied record does not specify the power-series variables, the precise target subgroup, or the construction of \(\mu\); these must be recovered from the cited sources before a fully self-contained definition can be given. The supplied archive reports that \(\mu\) is faithful, attributing this to A. Jaikin-Zapirain (2024), without independent verification here.
\end{description}

As additional background, the source reports that the Magnus homomorphism is injective on any subgroup of \(F^{\mathbb Q}\) of the indicated type
\[
\langle F,t\mid u=t^n\rangle,
\]
with attribution to G. Baumslag (1968). The supplied record does not explicitly specify the allowable parameters \(u,n\) or the embeddings identifying these presentations with subgroups of \(F^{\mathbb Q}\). This auxiliary assertion is retained with that qualification.
\subsection{Short English statement}
For an associative unital ring \(A\) with torsion-free additive group, the problem asks whether the source-defined free exponential group \(F^A\) is linear, has decidable universal theory, and can be characterized by length functions; it also asks whether free \(\mathbb Q\)-groups act freely on suitable \(\Lambda\)-trees. The supplied record labels the problem partially solved and reports affirmative solutions to the archived residual torsion-free nilpotence and Magnus-faithfulness questions. Several definitions and effectiveness conventions require clarification from the cited sources.
\subsection{Sources}
\begin{itemize}
\item Supplied local UnsolvedMath registry record: permanent ID \(700957\), exact source code \texttt{catalogue entry}, title \emph{Kourovka Notebook Problem 13.39}; upstream revision \texttt{d371ee917ac9647fbb06cd6fc529deec6f3dda82}. The supplied status is \texttt{partially\_solved}. Its status review is dated \texttt{2026-10-16} and reports parts (a) and (c) as archived solved, with the remaining parts in the main section. This is supplied dated curation, not independently verified literature assessment.
\item \emph{The Kourovka Notebook (Unsolved Problems in Group Theory)}, problem 13.39, originally in the 13th issue (1995), supplied as updated in the 21st issue (2026), identifier \texttt{arXiv:1401.0300v48}. The supplied collection credits E. I. Khukhro and V. D. Mazurov as editors. Proposed by A. G. Myasnikov and V. N. Remeslennikov; parts (a), (c), (e), and (f) are attributed to G. Baumslag.
\item A. G. Myasnikov and V. N. Remeslennikov, \emph{Siberian Math. J.} \textbf{35}, no. 5 (1994), 986--996. Supplied reference for the group-theoretic tensor product and free exponential group. No article title or URL is supplied.
\item A. G. Myasnikov and V. N. Remeslennikov, \emph{Int. J. Algebra Comput.} \textbf{6} (1996), 687--711. Supplied reference for a construction using free products with amalgamation. No article title or URL is supplied.
\item A. M. Gaglione, A. G. Myasnikov, V. N. Remeslennikov, and D. Spellman, \emph{Commun. Algebra} \textbf{25} (1997), 631--648. Supplied reference for the affirmative answer to part (b) when \(\langle1_A\rangle\) is pure in \((A,+)\). No article title or URL is supplied.
\item G. Baumslag, \emph{Commun. Pure Appl. Math.} \textbf{21} (1968), 491--506. Supplied reference for injectivity of the Magnus homomorphism on the indicated root-adjoining subgroups. No article title or URL is supplied.
\item R. Alperin and H. Bass, in \emph{Combinatorial group theory and topology, Alta, Utah, 1984}, Ann. Math. Stud. \textbf{111}, Princeton Univ. Press, 1987, 265--378. Supplied reference for the definition of a \(\Lambda\)-tree. No contribution title or URL is supplied.
\item A. Jaikin-Zapirain, \emph{Ann. Sci. \'Ec. Norm. Sup\'er. (4)} \textbf{57}, no. 4 (2024), 1101--1133. The supplied solved archive cites this reference for affirmative answers to parts (a) and (c). The full citation is present in the background, although the separate status-review excerpt for part (c) is truncated. No article title or URL is supplied. None of the supplied references has been independently consulted for this draft.
\end{itemize}
\end{document}
